\textbf{Phần 1. Trắc nghiệm lựa chọn}
\setcounterpageref{ex}{0}

\begin{ex} %[3D2V39]
Số lớn là 48, số bé là 6. Số lớn gấp mấy lần số bé?
\choice
{$7$ lần}
{$9$ lần}
{\True $8$ lần}
{$6$ lần}
\loigiai{
Số lớn gấp số bé số lần là: $48 : 6 = 8$ (lần).
\textbf{Đáp án của câu hỏi là C}
}
\end{ex}

\begin{ex} %[3D2V39]
Muốn tìm số lớn gấp mấy lần số bé, ta làm như sau:
\choice
{\True Lấy số lớn chia cho số bé}
{Lấy số lớn trừ đi số bé}
{Lấy số lớn nhân với số bé}
{Lấy số lớn cộng với số bé}
\loigiai{
Quy tắc toán học: Muốn biết số lớn gấp mấy lần số bé, ta lấy số lớn chia cho số bé.
\textbf{Đáp án của câu hỏi là A}
}
\end{ex}

\begin{ex} %[3D2V39]
Đoạn thẳng $AB$ dài $36\text{ cm}$, đoạn thẳng $CD$ dài $4\text{ cm}$. Hỏi đoạn thẳng $AB$ dài gấp mấy lần đoạn thẳng $CD$?
\choice
{$8$ lần}
{$7$ lần}
{$10$ lần}
{\True $9$ lần}
\loigiai{
$36 : 4 = 9$ (lần).
\textbf{Đáp án của câu hỏi là D}
}
\end{ex}

\begin{ex} %[3D2V39]
Năm nay bố 42 tuổi, con 7 tuổi. Hỏi năm nay tuổi bố gấp mấy lần tuổi con?
\choice
{$5$ lần}
{\True $6$ lần}
{$7$ lần}
{$8$ lần}
\loigiai{
Tuổi bố gấp tuổi con số lần là: $42 : 7 = 6$ (lần).
\textbf{Đáp án của câu hỏi là B}
}
\end{ex}

\begin{ex} %[3D2V39]
Một cuộn dây dài $60\text{ m}$. Người ta cắt lấy một đoạn dài $12\text{ m}$. Cuộn dây ban đầu dài gấp mấy lần đoạn dây đã cắt?
\choice
{$4$ lần}
{$6$ lần}
{\True $5$ lần}
{$3$ lần}
\loigiai{
$60 : 12 = 5$ (lần).
\textbf{Đáp án của câu hỏi là C}
}
\end{ex}

\begin{ex} %[3D2V39]
Trong vườn có 35 cây chuối và 7 cây bưởi. Số cây bưởi bằng một phần mấy số cây chuối?
\choice
{\True Một phần năm ($\dfrac{1}{5}$)}
{Một phần tư ($\dfrac{1}{4}$)}
{Một phần sáu ($\dfrac{1}{6}$)}
{Một phần bảy ($\dfrac{1}{7}$)}
\loigiai{
Số cây chuối gấp số cây bưởi: $35 : 7 = 5$ (lần). Vậy số cây bưởi bằng $\dfrac{1}{5}$ số cây chuối.
\textbf{Đáp án của câu hỏi là A}
}
\end{ex}

\begin{ex} %[3D2V39]
Số lớn là 72, số bé là 8. Số lớn nhiều hơn số bé bao nhiêu đơn vị và gấp số bé bao nhiêu lần?
\choice
{Nhiều hơn $64$ đơn vị và gấp $8$ lần}
{Nhiều hơn $56$ đơn vị và gấp $9$ lần}
{Nhiều hơn $70$ đơn vị và gấp $8$ lần}
{\True Nhiều hơn $64$ đơn vị và gấp $9$ lần}
\loigiai{
Nhiều hơn: $72 - 8 = 64$ đơn vị. Gấp: $72 : 8 = 9$ lần.
\textbf{Đáp án của câu hỏi là D}
}
\end{ex}

\begin{ex} %[3D2V39]
Một hình chữ nhật có chu vi $48\text{ cm}$, chiều rộng là $6\text{ cm}$. Chiều dài gấp mấy lần chiều rộng?
\choice
{$4$ lần}
{\True $3$ lần}
{$2$ lần}
{$5$ lần}
\loigiai{
Nửa chu vi: $48 : 2 = 24\text{ cm}$. Chiều dài: $24 - 6 = 18\text{ cm}$. Chiều dài gấp chiều rộng: $18 : 6 = 3$ (lần).
\textbf{Đáp án của câu hỏi là B}
}
\end{ex}

\begin{ex} %[3D2V39]
Hiện nay mẹ 36 tuổi, con 9 tuổi. Hỏi cách đây 3 năm, tuổi mẹ gấp mấy lần tuổi con?
\choice
{$4$ lần}
{$6$ lần}
{\True $5$ lần}
{$3$ lần}
\loigiai{
Cách đây 3 năm: tuổi mẹ là $36 - 3 = 33$ tuổi, tuổi con là $9 - 3 = 6$ tuổi (hoặc nếu xét khi con 6 tuổi mẹ 30 tuổi thì $30 : 6 = 5$ lần: hiện nay mẹ 33 tuổi con 9 tuổi; 3 năm trước mẹ 30, con 6; $30 : 6 = 5$ lần).
\textbf{Đáp án của câu hỏi là C}
}
\end{ex}

\begin{ex} %[3D2V39]
Có 40 con gà và 8 con vịt. Muốn số gà gấp 4 lần số vịt thì phải thêm vào bao nhiêu con vịt (giữ nguyên số gà)?
\choice
{\True $2$ con vịt}
{$4$ con vịt}
{$1$ con vịt}
{$3$ con vịt}
\loigiai{
Để số gà gấp 4 lần số vịt thì số vịt cần là: $40 : 4 = 10$ con. Cần thêm: $10 - 8 = 2$ con vịt.
\textbf{Đáp án của câu hỏi là A}
}
\end{ex}

\begin{ex} %[3D2V39]
Một can dầu chứa $45\ell$, một chai dầu chứa $5\ell$. Hỏi can dầu đựng gấp mấy lần chai dầu?
\choice
{$8$ lần}
{$7$ lần}
{$10$ lần}
{\True $9$ lần}
\loigiai{
$45 : 5 = 9$ (lần).
\textbf{Đáp án của câu hỏi là D}
}
\end{ex}

\begin{ex} %[3D2V39]
Số thứ nhất là 56. Số thứ hai kém số thứ nhất 49 đơn vị. Số thứ nhất gấp số thứ hai bao nhiêu lần?
\choice
{$7$ lần}
{\True $8$ lần}
{$6$ lần}
{$9$ lần}
\loigiai{
Số thứ hai là: $56 - 49 = 7$. Số thứ nhất gấp số thứ hai: $56 : 7 = 8$ (lần).
\textbf{Đáp án của câu hỏi là B}
}
\end{ex}

\textbf{Phần 2. Trắc nghiệm đúng sai}
\setcounterpageref{ex}{0}

\begin{ex} %[3D2V39]
Xét tính đúng sai của các câu nói về so sánh số lớn gấp mấy lần số bé:
\choiceTF
{\True Muốn biết số lớn gấp mấy lần số bé ta lấy số lớn chia cho số bé.}
{Muốn biết số lớn nhiều hơn số bé bao nhiêu đơn vị ta lấy số lớn chia cho số bé.}
{\True Nếu số lớn gấp 4 lần số bé thì số bé bằng một phần tư ($\dfrac{1}{4}$) số lớn.}
{Số 30 gấp 5 lần số 6, do đó số 30 nhiều hơn số 6 là 5 đơn vị.}
\loigiai{
\begin{itemize}
    \item Ý a) Đúng: Quy tắc chuẩn.
    \item Ý b) Sai: Sai vì muốn biết nhiều hơn bao nhiêu đơn vị ta dùng phép trừ.
    \item Ý c) Đúng: Mối quan hệ tỉ lệ thuận đảo.
    \item Ý d) Sai: Sai vì $30 - 6 = 24$ đơn vị.
\end{itemize}
\begin{itemize}
    \item a) \textbf{Đúng}.
    \item b) \textbf{Sai}.
    \item c) \textbf{Đúng}.
    \item d) \textbf{Sai}.
\end{itemize}
}
\end{ex}

\begin{ex} %[3D2V39]
Lớp 3A có 32 học sinh, trong đó có 8 học sinh giỏi:
\choiceTF
{\True Số học sinh cả lớp gấp 4 lần số học sinh giỏi.}
{\True Số học sinh không đạt loại giỏi là 24 bạn.}
{Số học sinh giỏi bằng $\dfrac{1}{3}$ số học sinh cả lớp.}
{\True Số học sinh không đạt loại giỏi gấp 3 lần số học sinh giỏi.}
\loigiai{
\begin{itemize}
    \item Ý a) Đúng: $32 : 8 = 4$ lần.
    \item Ý b) Đúng: $32 - 8 = 24$ bạn.
    \item Ý c) Sai: Sai vì bằng $\dfrac{1}{4}$ số học sinh cả lớp.
    \item Ý d) Đúng: $24 : 8 = 3$ lần.
\end{itemize}
\begin{itemize}
    \item a) \textbf{Đúng}.
    \item b) \textbf{Đúng}.
    \item c) \textbf{Sai}.
    \item d) \textbf{Đúng}.
\end{itemize}
}
\end{ex}

\textbf{Phần 3. Trả lời ngắn}
\setcounterpageref{ex}{0}

\begin{ex} %[3D2V39]
Số lớn là 72, số bé là 9. Hỏi số lớn gấp mấy lần số bé?
\shortans{8}
\loigiai{
\textbf{Đáp số:} $8$\\
$72 : 9 = 8$ lần.
}
\end{ex}

\begin{ex} %[3D2V39]
Một cây cau cao $18\text{ m}$, một cây ổi cao $3\text{ m}$. Hỏi cây cau cao gấp mấy lần cây ổi?
\shortans{6}
\loigiai{
\textbf{Đáp số:} $6$\\
$18 : 3 = 6$ lần.
}
\end{ex}

\begin{ex} %[3D2V39]
Bao ngô nặng $64\text{ kg}$, bao đậu nặng $8\text{ kg}$. Hỏi bao ngô nặng gấp mấy lần bao đậu?
\shortans{8}
\loigiai{
\textbf{Đáp số:} $8$\\
$64 : 8 = 8$ lần.
}
\end{ex}

\begin{ex} %[3D2V39]
Năm nay mẹ 35 tuổi, con 7 tuổi. Hỏi năm nay tuổi mẹ gấp mấy lần tuổi con?
\shortans{5}
\loigiai{
\textbf{Đáp số:} $5$\\
$35 : 7 = 5$ lần.
}
\end{ex}

\begin{ex} %[3D2V39]
Có 45 bông hoa hồng và 5 bông hoa cúc. Hỏi số hoa hồng gấp mấy lần số hoa cúc?
\shortans{9}
\loigiai{
\textbf{Đáp số:} $9$\\
$45 : 5 = 9$ lần.
}
\end{ex}

\begin{ex} %[3D2V39]
Thùng to có $54\ell$ nước, thùng nhỏ có $6\ell$ nước. Hỏi lượng nước thùng to gấp mấy lần thùng nhỏ?
\shortans{9}
\loigiai{
\textbf{Đáp số:} $9$\\
$54 : 6 = 9$ lần.
}
\end{ex}
